% ĐỢT 4: CỤM 3 (ĐỀ 1 & ĐỀ 2) - BỘ ĐỀ THAM KHẢO TUYỂN SINH 10 NĂM HỌC 2026-2027 SỞ GD&ĐT TP.HCM

% ==============================================================================
% CỤM 3 - ĐỀ THAM KHẢO 1 (Trang 15-16)
% ==============================================================================

\begin{bt}%[9D6H18-2]
% nguon: Đề tham khảo Tuyển sinh 10 năm học 2026-2027 Sở GD&ĐT TP.HCM - Cụm 3 Đề 1 - Bài 1
Cho hàm số $y = -\dfrac{1}{2}x^2$ có đồ thị là $(P)$.
\begin{enumerate}[a)]
	\item Vẽ đồ thị $(P)$ trên mặt phẳng tọa độ $Oxy$.
	\item Tìm tọa độ điểm $D$ thuộc $(P)$ ($D$ khác gốc tọa độ $O$) sao cho tung độ của điểm $D$ gấp đôi hoành độ của nó.
\end{enumerate}
\loigiai{
\begin{enumerate}[a)]
	\item Bảng giá trị của hàm số $y = -\dfrac{1}{2}x^2$:
	\begin{center}
		\begin{tabular}{|c|c|c|c|c|c|}
			\hline
			$x$ & $-4$ & $-2$ & $0$ & $2$ & $4$ \\ \hline
			$y$ & $-8$ & $-2$ & $0$ & $-2$ & $-8$ \\ \hline
		\end{tabular}
	\end{center}
	Đồ thị $(P)$ là một đường parabol đi qua các điểm $(-4; -8)$, $(-2; -2)$, $(0; 0)$, $(2; -2)$, $(4; -8)$, nhận trục $Oy$ làm trục đối xứng và có đỉnh là gốc tọa độ $O(0;0)$.
	\begin{center}
		\begin{tikzpicture}[>=stealth,scale=0.7]
			\draw[->] (-5,0) -- (5,0) node[below] {$x$};
			\draw[->] (0,-9) -- (0,1.5) node[left] {$y$};
			\draw (0,0) node[above right] {$O$};
			\foreach \x in {-4,-2,2,4}
				\draw (\x,0.1) -- (\x,-0.1) node[above] {\small $\x$};
			\foreach \y in {-8,-2}
				\draw (0.1,\y) -- (-0.1,\y) node[left] {\small $\y$};
			\draw[dashed] (-2,0) -- (-2,-2) -- (2,-2) -- (2,0);
			\draw[dashed] (-4,0) -- (-4,-8) -- (4,-8) -- (4,0);
			\draw[thick,smooth,samples=100,domain=-4.2:4.2] plot(\x,{-0.5*(\x)^2});
			\node[below right] at (2,-2) {$(P)$};
		\end{tikzpicture}
	\end{center}
	\item Gọi tọa độ điểm $D$ là $(x_D; y_D)$. Vì $D \in (P)$ nên $y_D = -\dfrac{1}{2}x_D^2$.\\
	Theo giả thiết, tung độ của $D$ gấp đôi hoành độ nên $y_D = 2x_D$.\\
	Ta có phương trình hoành độ giao điểm:
	$$-\dfrac{1}{2}x_D^2 = 2x_D \iff \dfrac{1}{2}x_D^2 + 2x_D = 0 \iff x_D\left(\dfrac{1}{2}x_D + 2\right) = 0 \iff \left[\begin{aligned} &x_D = 0 \\ &x_D = -4. \end{aligned}\right.$$
	Vì $D$ khác gốc tọa độ $O(0;0)$ nên $x_D = -4 \implies y_D = 2 \cdot (-4) = -8$.\\
	Vậy tọa độ điểm cần tìm là $D(-4; -8)$.
\end{enumerate}
}
\end{bt}

\begin{bt}%[9D6H20-2]
% nguon: Đề tham khảo Tuyển sinh 10 năm học 2026-2027 Sở GD&ĐT TP.HCM - Cụm 3 Đề 1 - Bài 2
Cho phương trình $x^2 - 4x - 2 = 0$.
\begin{enumerate}[a)]
	\item Chứng minh phương trình luôn có hai nghiệm phân biệt $x_1, x_2$.
	\item Không giải phương trình, hãy tính giá trị của biểu thức $A = \dfrac{x_1^2}{x_2} + \dfrac{x_2^2}{x_1}$.
\end{enumerate}
\loigiai{
\begin{enumerate}[a)]
	\item Phương trình $x^2 - 4x - 2 = 0$ có các hệ số $a = 1$, $b = -4$, $c = -2$.\\
	Ta có $\Delta' = (-2)^2 - 1 \cdot (-2) = 4 + 2 = 6 > 0$ (hoặc $a \cdot c = 1 \cdot (-2) = -2 < 0$).\\
	Do đó phương trình luôn có hai nghiệm phân biệt $x_1, x_2$.
	\item Theo định lí Vi-ét, ta có:
	$$\left\{\begin{aligned} &x_1 + x_2 = 4 \\ &x_1 x_2 = -2. \end{aligned}\right.$$
	Ta có:
	$$x_1^2 + x_2^2 = (x_1 + x_2)^2 - 2x_1 x_2 = 4^2 - 2 \cdot (-2) = 16 + 4 = 20.$$
	$$x_1^3 + x_2^3 = (x_1 + x_2)(x_1^2 - x_1 x_2 + x_2^2) = 4 \cdot [20 - (-2)] = 4 \cdot 22 = 88.$$
	Biểu thức:
	$$A = \dfrac{x_1^2}{x_2} + \dfrac{x_2^2}{x_1} = \dfrac{x_1^3 + x_2^3}{x_1 x_2} = \dfrac{88}{-2} = -44.$$
	Vậy $A = -44$.
\end{enumerate}
}
\end{bt}

\begin{bt}%[9D8H26-3]
% nguon: Đề tham khảo Tuyển sinh 10 năm học 2026-2027 Sở GD&ĐT TP.HCM - Cụm 3 Đề 1 - Bài 3
Một nhóm gồm $30$ học sinh lớp 9 được khảo sát về số giờ tự học mỗi tuần. Kết quả khảo sát được ghi lại như sau: có $4$ học sinh học $5$ giờ; $8$ học sinh học $6$ giờ; $10$ học sinh học $7$ giờ; $6$ học sinh học $8$ giờ; $2$ học sinh học $9$ giờ.
\begin{enumerate}[a)]
	\item Tính số giờ tự học trung bình mỗi tuần của nhóm học sinh trên.
	\item Chọn ngẫu nhiên một học sinh trong nhóm. Tính xác suất để học sinh được chọn có số giờ tự học mỗi tuần từ $8$ giờ trở lên.
	\item Chọn ngẫu nhiên một học sinh trong nhóm. Tính xác suất để học sinh được chọn có số giờ tự học mỗi tuần ít hơn $7$ giờ.
\end{enumerate}
\loigiai{
\begin{enumerate}[a)]
	\item Tổng số giờ tự học của cả nhóm $30$ học sinh là:
	$$4 \cdot 5 + 8 \cdot 6 + 10 \cdot 7 + 6 \cdot 8 + 2 \cdot 9 = 20 + 48 + 70 + 48 + 18 = 204\text{ (giờ)}.$$
	Số giờ tự học trung bình mỗi tuần của nhóm học sinh là:
	$$\bar{x} = \dfrac{204}{30} = 6{,}8\text{ (giờ)}.$$
	\item Gọi $M$ là biến cố: ``Học sinh được chọn có số giờ tự học mỗi tuần từ $8$ giờ trở lên''.\\
	Số học sinh có số giờ tự học từ $8$ giờ trở lên ($8$ giờ hoặc $9$ giờ) là: $6 + 2 = 8$ (học sinh).\\
	Xác suất của biến cố $M$ là:
	$$P(M) = \dfrac{8}{30} = \dfrac{4}{15}.$$
	\item Gọi $N$ là biến cố: ``Học sinh được chọn có số giờ tự học mỗi tuần ít hơn $7$ giờ''.\\
	Số học sinh có số giờ tự học ít hơn $7$ giờ ($5$ giờ hoặc $6$ giờ) là: $4 + 8 = 12$ (học sinh).\\
	Xác suất của biến cố $N$ là:
	$$P(N) = \dfrac{12}{30} = \dfrac{2}{5}.$$
\end{enumerate}
}
\end{bt}

\begin{bt}%[9D6V21-3]
% nguon: Đề tham khảo Tuyển sinh 10 năm học 2026-2027 Sở GD&ĐT TP.HCM - Cụm 3 Đề 1 - Bài 4
Một sân vận động có thiết kế mặt sân gồm một hình chữ nhật kích thước $x$ mét và $y$ mét, cùng hai đầu là hai nửa hình tròn có đường kính bằng $x$ mét (với $x, y > 0$).
\begin{enumerate}[a)]
	\item Hãy viết biểu thức biểu diễn chu vi $P$ của toàn bộ sân vận động theo $x, y$ và số $\pi$.
	\item Cho biết chiều dài hình chữ nhật gấp rưỡi chiều rộng ($y = 1{,}5x$), chu vi toàn bộ sân vận động là $314{,}16\text{ m}$. Lấy $\pi \approx 3{,}14$, hãy tính diện tích mặt sân và chi phí trải thảm cỏ nhân tạo cho toàn bộ mặt sân biết giá mỗi mét vuông cỏ là $250\,000$ đồng (kết quả các độ dài và diện tích làm tròn đến chữ số thập phân thứ hai, chi phí làm tròn đến hàng đơn vị).
\end{enumerate}
\loigiai{
\begin{enumerate}[a)]
	\item Chu vi của sân vận động gồm độ dài hai cạnh của hình chữ nhật có chiều dài $y$ và chu vi của hai nửa hình tròn ghép lại thành một hình tròn đường kính $x$ (chu vi bằng $\pi x$).\\
	Biểu thức biểu diễn chu vi của toàn bộ sân vận động là:
	$$P = 2y + \pi x\text{ (m)}.$$
	\item Theo đề bài, $y = 1{,}5x$ và $P = 314{,}16\text{ m}$. Thay vào biểu thức chu vi, ta được:
	$$2(1{,}5x) + 3{,}14x = 314{,}16 \iff 3x + 3{,}14x = 314{,}16 \iff 6{,}14x = 314{,}16 \implies x = \dfrac{314{,}16}{6{,}14} \approx 51{,}17\text{ (m)}.$$
	Chiều dài của hình chữ nhật là:
	$$y = 1{,}5x \approx 1{,}5 \cdot 51{,}17 = 76{,}75\text{ (m)}.$$
	Diện tích toàn bộ sân vận động gồm diện tích hình chữ nhật và diện tích hình tròn có bán kính $r = \dfrac{x}{2} \approx \dfrac{51{,}17}{2} = 25{,}585\text{ (m)}$:
	$$S = x \cdot y + \pi \cdot \left(\dfrac{x}{2}\right)^2 \approx 51{,}17 \cdot 76{,}75 + 3{,}14 \cdot (25{,}585)^2 \approx 3927{,}2975 + 2055{,}4285 \approx 5982{,}07\text{ (m}^2\text{)}.$$
	Chi phí trải thảm cỏ nhân tạo cho toàn bộ mặt sân là:
	$$5982{,}07 \cdot 250\,000 = 1\,495\,517\,500\text{ (đồng)}.$$
	(Nếu lấy giá trị chính xác $x = \dfrac{314{,}16}{6{,}14}$ thì chi phí xấp xỉ $1\,495\,516\,632$ đồng).
\end{enumerate}
}
\end{bt}

\begin{bt}%[9H0V31-3]
% nguon: Đề tham khảo Tuyển sinh 10 năm học 2026-2027 Sở GD&ĐT TP.HCM - Cụm 3 Đề 1 - Bài 5
Một cơ sở sản xuất nến thơm dạng hình cầu có bán kính $R = 4\text{ cm}$. Để đóng gói sản phẩm, người ta đặt nến thơm vào một hộp nhựa trong suốt dạng hình trụ có kích thước vừa khít với quả cầu nến (đường kính đáy bằng đường kính quả cầu và chiều cao hình trụ bằng đường kính quả cầu). Lấy $\pi \approx 3{,}14$.
\begin{enumerate}[a)]
	\item Tính thể tích của hộp nhựa hình trụ và thể tích của quả cầu nến thơm (làm tròn kết quả đến chữ số thập phân thứ hai).
	\item Người ta đổ thêm các hạt xốp mềm vào phần không gian còn trống giữa quả cầu nến và vỏ hộp hình trụ để chống sốc. Biết chi phí mua hạt xốp là $150$ đồng cho mỗi centimét khối hạt xốp. Tính số tiền mua hạt xốp để đóng gói $500$ bộ nến thơm như trên (kết quả làm tròn đến hàng nghìn).
\end{enumerate}
\loigiai{
\begin{enumerate}[a)]
	\item Quả cầu nến có bán kính $R = 4\text{ cm}$.\\
	Hộp hình trụ có bán kính đáy $r = R = 4\text{ cm}$ và chiều cao $h = 2R = 8\text{ cm}$.\\
	Thể tích của hộp nhựa hình trụ là:
	$$V_{\text{trụ}} = \pi r^2 h \approx 3{,}14 \cdot 4^2 \cdot 8 = 401{,}92\text{ (cm}^3\text{)}.$$
	Thể tích của quả cầu nến là:
	$$V_{\text{cầu}} = \dfrac{4}{3}\pi R^3 \approx \dfrac{4}{3} \cdot 3{,}14 \cdot 4^3 \approx 267{,}95\text{ (cm}^3\text{)}.$$
	\item Thể tích phần không gian còn trống trong mỗi hộp là:
	$$V_{\text{trống}} = V_{\text{trụ}} - V_{\text{cầu}} \approx 401{,}92 - 267{,}95 = 133{,}97\text{ (cm}^3\text{)}.$$
	Tổng thể tích hạt xốp cần dùng để đóng gói $500$ hộp là:
	$$V_{\text{tổng}} = 500 \cdot 133{,}97 = 66\,985\text{ (cm}^3\text{)}.$$
	Số tiền mua hạt xốp cho $500$ bộ sản phẩm là:
	$$66\,985 \cdot 150 = 10\,047\,750\text{ (đồng)}.$$
	Làm tròn đến hàng nghìn là $10\,048\,000$ đồng.
\end{enumerate}
}
\end{bt}

\begin{bt}%[9D1H3-4]
% nguon: Đề tham khảo Tuyển sinh 10 năm học 2026-2027 Sở GD&ĐT TP.HCM - Cụm 3 Đề 1 - Bài 6
Hai người đi xe máy cùng khởi hành một lúc từ địa điểm $A$ để đi đến địa điểm $B$. Người thứ nhất đi với vận tốc ban đầu lớn hơn vận tốc của người thứ hai là $20\%$. Sau khi đi được $30$ phút, do gặp đoạn đường xấu, người thứ nhất buộc phải giảm vận tốc xuống chỉ còn bằng một nửa vận tốc của người thứ hai trên quãng đường còn lại. Biết rằng cả hai người đến địa điểm $B$ cùng một lúc. Hỏi thời gian hai người đi từ $A$ đến $B$ là bao nhiêu phút?
\loigiai{
Đổi: $30\text{ phút} = 0{,}5\text{ giờ}$.\\
Gọi $v_2$ (km/h) là vận tốc của người thứ hai ($v_2 > 0$).\\
Gọi $t$ (giờ) là tổng thời gian hai người đi hết quãng đường $AB$ ($t > 0{,}5$).\\
Vận tốc ban đầu của người thứ nhất là:
$$v_1 = v_2 + 20\% v_2 = 1{,}2 v_2\text{ (km/h)}.$$
Vận tốc của người thứ nhất trên đoạn đường còn lại là:
$$v'_1 = 0{,}5 v_2\text{ (km/h)}.$$
Quãng đường người thứ hai đi được trong thời gian $t$ là:
$$S_{AB} = v_2 \cdot t\text{ (km)}.$$
Quãng đường người thứ nhất đi được gồm hai đoạn:
\begin{itemize}
	\item Đoạn đầu đi trong $0{,}5$ giờ: $S_1 = 1{,}2 v_2 \cdot 0{,}5 = 0{,}6 v_2\text{ (km)}$.
	\item Đoạn sau đi trong $t - 0{,}5$ giờ: $S_2 = 0{,}5 v_2 \cdot (t - 0{,}5)\text{ (km)}$.
\end{itemize}
Vì cả hai cùng đi hết quãng đường $AB$ và đến nơi cùng lúc, ta có phương trình:
$$v_2 \cdot t = 0{,}6 v_2 + 0{,}5 v_2 (t - 0{,}5).$$
Chia cả hai vế cho $v_2 > 0$, ta được:
$$t = 0{,}6 + 0{,}5(t - 0{,}5) \iff t = 0{,}6 + 0{,}5t - 0{,}25 \iff 0{,}5t = 0{,}35 \iff t = 0{,}7\text{ (giờ)}.$$
Đổi $0{,}7\text{ giờ} = 0{,}7 \cdot 60 = 42\text{ phút}$.\\
Vậy thời gian hai người đi từ $A$ đến $B$ là $42$ phút.
}
\end{bt}

\begin{bt}%[9H9C29-3]
% nguon: Đề tham khảo Tuyển sinh 10 năm học 2026-2027 Sở GD&ĐT TP.HCM - Cụm 3 Đề 1 - Bài 7
Cho tam giác nhọn $ABC$ ($AB < AC$) nội tiếp đường tròn $(O; R)$ với $R = 5\text{ cm}$. Các đường cao $AD$ và $BE$ của tam giác $ABC$ cắt nhau tại trực tâm $H$. Kẻ đường kính $AK$ của đường tròn $(O)$.
\begin{enumerate}[a)]
	\item Chứng minh $AC \perp KC$ và tam giác $ABD$ đồng dạng với tam giác $AKC$.
	\item Chứng minh tứ giác $CDHE$ nội tiếp và $AH \cdot AD = AE \cdot AC$.
	\item Cho biết độ dài cạnh $BC = 8\text{ cm}$. Tính độ dài đoạn thẳng $AH$.
\end{enumerate}
\loigiai{
\begin{center}
	\begin{tikzpicture}[scale=0.9,>=stealth]
		\coordinate (O) at (0,0);
		\def\R{3.5}
		\draw[thick] (O) circle (\R);
		\coordinate (A) at (-0.8,{\R*sin(acos(-0.8/\R))});
		\coordinate (B) at ({-3.5*cos(25)},{-3.5*sin(25)});
		\coordinate (C) at ({3.5*cos(15)},{-3.5*sin(15)});
		\coordinate (K) at ($(O)!-1!(A)$);
		% Duong cao AD
		\coordinate (D) at ($(B)!(A)!(C)$);
		% Duong cao BE
		\coordinate (E) at ($(A)!(B)!(C)$);
		\coordinate (H) at (intersection of A--D and B--E);
		\coordinate (M) at ($(B)!0.5!(C)$);

		\draw[thick] (A) -- (B) -- (C) -- cycle;
		\draw[thick] (A) -- (D);
		\draw[thick] (B) -- (E);
		\draw[thick] (A) -- (K) -- (C);
		\draw[dashed] (B) -- (K);
		\draw[dashed] (C) -- (K);
		\draw[dashed] (O) -- (M);
		\draw[dashed] (H) -- (C);

		\fill (A) circle (1.5pt) node[above] {$A$};
		\fill (B) circle (1.5pt) node[left] {$B$};
		\fill (C) circle (1.5pt) node[right] {$C$};
		\fill (O) circle (1.5pt) node[below right] {$O$};
		\fill (D) circle (1.5pt) node[below] {$D$};
		\fill (E) circle (1.5pt) node[above right] {$E$};
		\fill (H) circle (1.5pt) node[left] {$H$};
		\fill (K) circle (1.5pt) node[below left] {$K$};
		\fill (M) circle (1.5pt) node[below] {$M$};
	\end{tikzpicture}
\end{center}
\begin{enumerate}[a)]
	\item Ta có $\widehat{ACK} = 90^\circ$ (góc nội tiếp chắn nửa đường tròn $(O)$ đường kính $AK$).\\
	Suy ra $AC \perp KC$.\\
	Xét $\Delta ABD$ và $\Delta AKC$ có:
	\begin{itemize}
		\item $\widehat{ADB} = \widehat{ACK} = 90^\circ$;
		\item $\widehat{ABD} = \widehat{AKC}$ (hai góc nội tiếp cùng chắn cung $AC$ của $(O)$).
	\end{itemize}
	Do đó $\Delta ABD \backsim \Delta AKC$ (g-g).
	\item Vì $AD \perp BC$ tại $D$ nên $\widehat{HD C} = 90^\circ$.\\
	Vì $BE \perp AC$ tại $E$ nên $\widehat{HEC} = 90^\circ$.\\
	Xét tứ giác $CDHE$ có $\widehat{HDC} + \widehat{HEC} = 90^\circ + 90^\circ = 180^\circ$, suy ra tứ giác $CDHE$ nội tiếp.\\
	Xét $\Delta AHE$ và $\Delta ACD$ có:
	\begin{itemize}
		\item $\widehat{CAD}$ chung;
		\item $\widehat{AEH} = \widehat{ADC} = 90^\circ$.
	\end{itemize}
	Do đó $\Delta AHE \backsim \Delta ACD$ (g-g), suy ra $\dfrac{AH}{AC} = \dfrac{AE}{AD} \implies AH \cdot AD = AE \cdot AC$.
	\item Ta có $\widehat{ABK} = 90^\circ$ (góc nội tiếp chắn nửa đường tròn), nên $BK \perp AB$.\\
	Lại có $CH \perp AB$ (vì $H$ là trực tâm tam giác $ABC$). Do đó $BK \parallel CH$.\\
	Tương tự, $CK \perp AC$ và $BH \perp AC$ nên $CK \parallel BH$.\\
	Tứ giác $BHCK$ có hai cặp cạnh đối song song nên là hình bình hành.\\
	Gọi $M$ là trung điểm của $BC$, suy ra $M$ cũng là trung điểm của đường chéo $HK$.\\
	Xét tam giác $AHK$, có $O$ là trung điểm $AK$, $M$ là trung điểm $HK$, nên $OM$ là đường trung bình của tam giác $AHK$.\\
	Suy ra $AH = 2 \cdot OM$.\\
	Vì $M$ là trung điểm của dây $BC$ không đi qua tâm $O$ nên $OM \perp BC$ tại $M$.\\
	Do đó tam giác $OBM$ vuông tại $M$. Áp dụng định lí Pythagore trong $\Delta OBM$:
	$$OM = \sqrt{OB^2 - BM^2} = \sqrt{R^2 - \left(\dfrac{BC}{2}\right)^2} = \sqrt{5^2 - 4^2} = 3\text{ (cm)}.$$
	Vậy $AH = 2 \cdot OM = 2 \cdot 3 = 6\text{ (cm)}$.
\end{enumerate}
}
\end{bt}


% ==============================================================================
% CỤM 3 - ĐỀ THAM KHẢO 2 (Trang 17-18)
% ==============================================================================

\begin{bt}%[9D6H18-2]
% nguon: Đề tham khảo Tuyển sinh 10 năm học 2026-2027 Sở GD&ĐT TP.HCM - Cụm 3 Đề 2 - Bài 1
Cho hàm số $y = \dfrac{x^2}{2}$ có đồ thị là parabol $(P)$.
\begin{enumerate}[a)]
	\item Vẽ đồ thị parabol $(P)$ trên mặt phẳng tọa độ $Oxy$.
	\item Tìm các điểm thuộc đồ thị $(P)$ có tung độ bằng $4$.
\end{enumerate}
\loigiai{
\begin{enumerate}[a)]
	\item Bảng giá trị của hàm số $y = \dfrac{x^2}{2}$:
	\begin{center}
		\begin{tabular}{|c|c|c|c|c|c|}
			\hline
			$x$ & $-4$ & $-2$ & $0$ & $2$ & $4$ \\ \hline
			$y$ & $8$ & $2$ & $0$ & $2$ & $8$ \\ \hline
		\end{tabular}
	\end{center}
	Đồ thị $(P)$ là một parabol đi qua các điểm $(-4; 8)$, $(-2; 2)$, $(0; 0)$, $(2; 2)$, $(4; 8)$, có đỉnh là gốc tọa độ $O(0;0)$ và nhận trục tung $Oy$ làm trục đối xứng.
	\begin{center}
		\begin{tikzpicture}[>=stealth,scale=0.7]
			\draw[->] (-5,0) -- (5,0) node[below] {$x$};
			\draw[->] (0,-1) -- (0,9) node[left] {$y$};
			\draw (0,0) node[below right] {$O$};
			\foreach \x in {-4,-2,2,4}
				\draw (\x,0.1) -- (\x,-0.1) node[below] {\small $\x$};
			\foreach \y in {2,4,8}
				\draw (0.1,\y) -- (-0.1,\y) node[left] {\small $\y$};
			\draw[dashed] (-2,0) -- (-2,2) -- (2,2) -- (2,0);
			\draw[dashed] (-4,0) -- (-4,8) -- (4,8) -- (4,0);
			\draw[thick,smooth,samples=100,domain=-4.1:4.1] plot(\x,{0.5*(\x)^2});
			\node[above left] at (2,2) {$(P)$};
		\end{tikzpicture}
	\end{center}
	\item Gọi điểm cần tìm là $N(x; y)$. Vì điểm $N$ thuộc đồ thị $(P)$ và có tung độ bằng $4$ nên:
	$$4 = \dfrac{x^2}{2} \iff x^2 = 8 \iff x = \pm 2\sqrt{2}.$$
	Vậy có hai điểm thỏa mãn yêu cầu bài toán là $(2\sqrt{2}; 4)$ và $(-2\sqrt{2}; 4)$.
\end{enumerate}
}
\end{bt}

\begin{bt}%[9D6H20-2]
% nguon: Đề tham khảo Tuyển sinh 10 năm học 2026-2027 Sở GD&ĐT TP.HCM - Cụm 3 Đề 2 - Bài 2
Cho phương trình bậc hai $3x^2 + 2x - 3 = 0$.
\begin{enumerate}[a)]
	\item Chứng minh phương trình luôn có hai nghiệm phân biệt $x_1, x_2$.
	\item Không giải phương trình, hãy tính giá trị của biểu thức:
	$$M = (x_1 - 2x_2)(x_2 - x_1) + x_2^2.$$
\end{enumerate}
\loigiai{
\begin{enumerate}[a)]
	\item Phương trình $3x^2 + 2x - 3 = 0$ có các hệ số $a = 3$, $b = 2$, $c = -3$.\\
	Ta có $a \cdot c = 3 \cdot (-3) = -9 < 0$.\\
	Do $a$ và $c$ trái dấu nên phương trình luôn có hai nghiệm phân biệt $x_1, x_2$.
	\item Theo định lí Vi-ét, ta có:
	$$\left\{\begin{aligned} &x_1 + x_2 = -\dfrac{2}{3} \\ &x_1 x_2 = \dfrac{-3}{3} = -1. \end{aligned}\right.$$
	Biến đổi biểu thức $M$:
	\begin{align*}
		M &= (x_1 - 2x_2)(x_2 - x_1) + x_2^2 \\
		&= x_1 x_2 - x_1^2 - 2x_2^2 + 2x_1 x_2 + x_2^2 \\
		&= 3x_1 x_2 - (x_1^2 + x_2^2) \\
		&= 3x_1 x_2 - [(x_1 + x_2)^2 - 2x_1 x_2] \\
		&= 5x_1 x_2 - (x_1 + x_2)^2.
	\end{align*}
	Thay $x_1 + x_2 = -\dfrac{2}{3}$ và $x_1 x_2 = -1$ vào biểu thức $M$, ta được:
	$$M = 5 \cdot (-1) - \left(-\dfrac{2}{3}\right)^2 = -5 - \dfrac{4}{9} = -\dfrac{49}{9}.$$
	Vậy $M = -\dfrac{49}{9}$.
\end{enumerate}
}
\end{bt}

\begin{bt}%[9D8H26-3]
% nguon: Đề tham khảo Tuyển sinh 10 năm học 2026-2027 Sở GD&ĐT TP.HCM - Cụm 3 Đề 2 - Bài 3
Bảng thống kê sau đây ghi lại số liệu thu nhập và chi phí (đơn vị: triệu đồng) trong một tuần của một cửa hàng:
\begin{center}
	\begin{tabular}{|l|c|c|c|c|c|c|c|}
		\hline
		\textbf{Ngày trong tuần} & \textbf{Thứ 2} & \textbf{Thứ 3} & \textbf{Thứ 4} & \textbf{Thứ 5} & \textbf{Thứ 6} & \textbf{Thứ 7} & \textbf{Chủ nhật} \\ \hline
		\textbf{Thu nhập (triệu đồng)} & $150$ & $180$ & $160$ & $220$ & $210$ & $250$ & $280$ \\ \hline
		\textbf{Chi phí (triệu đồng)} & $120$ & $140$ & $130$ & $150$ & $170$ & $210$ & $200$ \\ \hline
	\end{tabular}
\end{center}
Biết lợi nhuận của cửa hàng được tính bằng thu nhập trừ đi chi phí.
\begin{enumerate}[a)]
	\item Ngày nào trong tuần cửa hàng đạt mức lợi nhuận cao nhất và lợi nhuận của ngày đó là bao nhiêu triệu đồng?
	\item Chọn ngẫu nhiên một ngày trong tuần.
	\begin{itemize}
		\item Tính xác suất để ngày được chọn có mức thu nhập lớn hơn $200$ triệu đồng.
		\item Tính xác suất để ngày được chọn có lợi nhuận lớn hơn $50$ triệu đồng.
	\end{itemize}
\end{enumerate}
\loigiai{
\begin{enumerate}[a)]
	\item Bảng tính lợi nhuận các ngày trong tuần của cửa hàng:
	\begin{center}
		\begin{tabular}{|l|c|c|c|c|c|c|c|}
			\hline
			\textbf{Ngày} & Thứ 2 & Thứ 3 & Thứ 4 & Thứ 5 & Thứ 6 & Thứ 7 & Chủ nhật \\ \hline
			\textbf{Lợi nhuận (triệu đồng)} & $30$ & $40$ & $30$ & $70$ & $40$ & $40$ & $80$ \\ \hline
		\end{tabular}
	\end{center}
	Cửa hàng đạt mức lợi nhuận cao nhất vào ngày Chủ nhật với mức lợi nhuận là:
	$$280 - 200 = 80\text{ (triệu đồng)}.$$
	\item Không gian mẫu gồm $7$ ngày trong tuần, các kết quả có cùng khả năng xảy ra nên $n(\Omega) = 7$.
	\begin{itemize}
		\item Gọi $A$ là biến cố: ``Ngày được chọn có thu nhập lớn hơn $200$ triệu đồng''.\\
		Các ngày có thu nhập lớn hơn $200$ triệu đồng là: Thứ 5 ($220$), Thứ 6 ($210$), Thứ 7 ($250$), Chủ nhật ($280$).\\
		Có $4$ ngày thỏa mãn, suy ra xác suất của biến cố $A$ là:
		$$P(A) = \dfrac{4}{7}.$$
		\item Gọi $B$ là biến cố: ``Ngày được chọn có lợi nhuận lớn hơn $50$ triệu đồng''.\\
		Các ngày có lợi nhuận lớn hơn $50$ triệu đồng là: Thứ 5 ($70$ triệu đồng) và Chủ nhật ($80$ triệu đồng).\\
		Có $2$ ngày thỏa mãn, suy ra xác suất của biến cố $B$ là:
		$$P(B) = \dfrac{2}{7}.$$
	\end{itemize}
\end{enumerate}
}
\end{bt}

\begin{bt}%[9D6V21-3]
% nguon: Đề tham khảo Tuyển sinh 10 năm học 2026-2027 Sở GD&ĐT TP.HCM - Cụm 3 Đề 2 - Bài 4
Một khu vườn hình vuông có độ dài cạnh là $3x$ mét ($x > 0$). Ở chính giữa khu vườn, người ta xây một hồ cá hình tròn nội tiếp tiếp xúc với bốn cạnh của hình vuông có đường kính bằng cạnh hình vuông (bán kính hồ là $r = \dfrac{3x}{2}$ mét). Toàn bộ phần đất còn lại trong vườn được dùng để trồng cỏ nhung Nhật Bản.
\begin{enumerate}[a)]
	\item Viết biểu thức biểu diễn diện tích của hồ cá theo $x$ và số $\pi$.
	\item Gia đình đã mua cỏ với số lượng dự phòng hao hụt bằng $10\%$ so với diện tích đất trồng cỏ thực tế. Tổng số tiền gia đình đã trả để mua cỏ là $425\,700$ đồng với giá $50\,000$ đồng cho mỗi mét vuông cỏ. Lấy $\pi \approx 3{,}14$, hãy tính độ dài cạnh của khu vườn và diện tích của toàn bộ khu vườn.
\end{enumerate}
\loigiai{
\begin{enumerate}[a)]
	\item Bán kính hồ cá hình tròn là $r = \dfrac{3x}{2}\text{ (m)}$.\\
	Biểu thức biểu diễn diện tích của hồ cá là:
	$$S_{\text{hồ}} = \pi r^2 = \pi \left(\dfrac{3x}{2}\right)^2 = \dfrac{9\pi}{4}x^2\text{ (m}^2\text{)}.$$
	\item Diện tích của khu vườn hình vuông là:
	$$S_{\text{vườn}} = (3x)^2 = 9x^2\text{ (m}^2\text{)}.$$
	Diện tích đất thực tế dùng để trồng cỏ là:
	$$S_{\text{cỏ}} = S_{\text{vườn}} - S_{\text{hồ}} = 9x^2 - \dfrac{9 \cdot 3{,}14}{4}x^2 = 9x^2 - 7{,}065x^2 = 1{,}935x^2\text{ (m}^2\text{)}.$$
	Diện tích cỏ đã mua tính cả $10\%$ hao hụt là:
	$$S_{\text{mua}} = S_{\text{cỏ}} + 10\% S_{\text{cỏ}} = 1{,}1 \cdot S_{\text{cỏ}}.$$
	Theo đề bài, số tiền mua cỏ là $425\,700$ đồng với đơn giá $50\,000$ đồng/$\text{m}^2$, nên diện tích cỏ đã mua là:
	$$S_{\text{mua}} = \dfrac{425\,700}{50\,000} = 8{,}514\text{ (m}^2\text{)}.$$
	Do đó diện tích trồng cỏ thực tế là:
	$$S_{\text{cỏ}} = \dfrac{8{,}514}{1{,}1} = 7{,}74\text{ (m}^2\text{)}.$$
	Ta có phương trình:
	$$1{,}935x^2 = 7{,}74 \iff x^2 = \dfrac{7{,}74}{1{,}935} = 4 \implies x = 2\text{ (m) (do } x > 0\text{)}.$$
	Độ dài cạnh của khu vườn là: $3x = 3 \cdot 2 = 6\text{ (m)}$.\\
	Diện tích của toàn bộ khu vườn là: $S_{\text{vườn}} = 6^2 = 36\text{ (m}^2\text{)}$.
\end{enumerate}
}
\end{bt}

\begin{bt}%[9H0V31-3]
% nguon: Đề tham khảo Tuyển sinh 10 năm học 2026-2027 Sở GD&ĐT TP.HCM - Cụm 3 Đề 2 - Bài 5
Một thùng tôn hình trụ có chiều cao $h = 36\text{ cm}$ và đường kính đáy là $3\text{ dm}$.
\begin{enumerate}[a)]
	\item Tính thể tích của thùng tôn hình trụ đó theo đơn vị lít (làm tròn kết quả đến chữ số thập phân thứ hai). Lấy $\pi \approx 3{,}14$.
	\item Người ta dùng thùng tôn trên để múc nước từ một giếng đổ vào một bể chứa nước rỗng có dung tích $1\text{ m}^3$. Mỗi lần múc, lượng nước trong thùng chỉ đạt $90\%$ thể tích tối đa của thùng. Hỏi sau $40$ lần múc đổ vào bể thì bể đã đầy nước chưa? Vì sao?
\end{enumerate}
\loigiai{
\begin{enumerate}[a)]
	\item Đổi: $3\text{ dm} = 30\text{ cm}$. Bán kính đáy của thùng tôn là $R = \dfrac{30}{2} = 15\text{ cm}$.\\
	Thể tích của thùng tôn hình trụ là:
	$$V = \pi R^2 h \approx 3{,}14 \cdot 15^2 \cdot 36 = 25\,434\text{ (cm}^3\text{)}.$$
	Đổi: $25\,434\text{ cm}^3 = 25{,}434\text{ dm}^3 \approx 25{,}43\text{ lít}$ (hoặc nếu dùng $\pi$ chính xác thì $V = 8100\pi\text{ cm}^3 \approx 25{,}45\text{ lít}$).
	\item Thể tích nước múc được trong mỗi lần là:
	$$V_1 = 90\% \cdot V = 0{,}9 \cdot 25\,434 = 22\,890{,}6\text{ (cm}^3\text{)}.$$
	Tổng thể tích nước đổ vào bể sau $40$ lần múc là:
	$$V_{40} = 40 \cdot V_1 = 40 \cdot 22\,890{,}6 = 915\,624\text{ (cm}^3\text{)}.$$
	Đổi sang đơn vị mét khối:
	$$915\,624\text{ cm}^3 = 0{,}915624\text{ m}^3.$$
	Vì $0{,}915624\text{ m}^3 < 1\text{ m}^3$ nên sau $40$ lần múc, lượng nước đổ vào vẫn chưa đủ để làm đầy bể nước.
\end{enumerate}
}
\end{bt}

\begin{bt}%[9D1H3-4]
% nguon: Đề tham khảo Tuyển sinh 10 năm học 2026-2027 Sở GD&ĐT TP.HCM - Cụm 3 Đề 2 - Bài 6
Một phân xưởng dự định giao cho hai máy tự động gia công một lô linh kiện. Mỗi giờ cả hai máy gia công được số lượng linh kiện như nhau. Nếu tăng năng suất của mỗi máy thêm $3$ linh kiện mỗi giờ thì thời gian hoàn thành lô hàng giảm đi $2$ giờ so với dự kiến. Nếu tăng năng suất của mỗi máy thêm $5$ linh kiện mỗi giờ thì trong thời gian dự kiến ban đầu giảm đi $3$ giờ, hai máy đã sản xuất vượt chỉ tiêu $15$ linh kiện. Hỏi lô linh kiện theo dự định gồm bao nhiêu chiếc?
\loigiai{
Gọi $x$ là số linh kiện mỗi giờ cả hai máy gia công được theo dự định ($x > 0$).\\
Gọi $y$ (giờ) là thời gian hai máy hoàn thành lô hàng theo kế hoạch dự định ($y > 3$).\\
Tổng số linh kiện của lô hàng theo dự định là $x \cdot y$ (linh kiện).\\
Vì cả hai máy có năng suất như nhau nên mỗi máy gia công được $\dfrac{x}{2}$ linh kiện mỗi giờ.
\begin{itemize}
	\item Khi tăng năng suất mỗi máy thêm $3$ linh kiện/giờ thì năng suất của cả hai máy tăng thêm $6$ linh kiện/giờ, tức là $x + 6$ linh kiện/giờ (hoặc nếu $x$ là năng suất của mỗi máy thì tổng sản phẩm là $2xy$).
\end{itemize}
Theo cách hiểu chuẩn của đề bài: Gọi $x$ là số linh kiện cả hai máy làm trong một giờ và $y$ là số giờ dự định.\\
Trường hợp 1: Nếu năng suất của hệ thống tăng thêm $3$ linh kiện/giờ và thời gian giảm $2$ giờ:
$$(x + 3)(y - 2) = xy \iff xy - 2x + 3y - 6 = xy \iff -2x + 3y = 6.$$
Trường hợp 2: Nếu tăng thêm $5$ linh kiện/giờ thì trong thời gian $y - 3$ giờ sản xuất được $xy + 15$ linh kiện:
$$(x + 5)(y - 3) = xy + 15 \iff xy - 3x + 5y - 15 = xy + 15 \iff -3x + 5y = 30.$$
Ta có hệ phương trình bậc nhất hai ẩn:
$$\left\{\begin{aligned} &-2x + 3y = 6 \\ &-3x + 5y = 30 \end{aligned}\right. \iff \left\{\begin{aligned} &-6x + 9y = 18 \\ &-6x + 10y = 60 \end{aligned}\right. \iff \left\{\begin{aligned} &y = 42 \\ &2x = 3 \cdot 42 - 6 \implies x = 60. \end{aligned}\right.$$
Cả hai giá trị $x = 60$ và $y = 42$ đều thỏa mãn điều kiện bài toán.\\
Vậy số linh kiện của lô hàng dự định sản xuất là:
$$60 \cdot 42 = 2520\text{ (linh kiện)}.$$
}
\end{bt}

\begin{bt}%[9H9C29-3]
% nguon: Đề tham khảo Tuyển sinh 10 năm học 2026-2027 Sở GD&ĐT TP.HCM - Cụm 3 Đề 2 - Bài 7
Cho đường tròn $(O; R)$ và điểm $M$ nằm bên ngoài đường tròn sao cho $OM = 2R$. Kẻ các tiếp tuyến $MA, MC$ với đường tròn $(O)$ ($A, C$ là các tiếp điểm). Kẻ đường kính $AB$ của $(O)$. Đoạn thẳng $MB$ cắt đường tròn $(O)$ tại điểm $D$ ($D$ khác $B$). Đoạn thẳng $OM$ cắt $AC$ tại $H$.
\begin{enumerate}[a)]
	\item Chứng minh tam giác $ABD$ vuông tại $D$ và $OM \perp AC$ tại $H$.
	\item Chứng minh $OD^2 = OH \cdot OM$ và $\widehat{ODH} = \widehat{DAC}$.
	\item Gọi $K$ là trung điểm của đoạn thẳng $BD$. Tia $AC$ cắt tia $OK$ tại $E$. Chứng minh $OE \perp AB$ và tính diện tích tam giác $OEB$ theo $R$.
\end{enumerate}
\loigiai{
\begin{center}
	\begin{tikzpicture}[scale=0.9,>=stealth]
		\coordinate (O) at (0,0);
		\def\R{2.8}
		\draw[thick] (O) circle (\R);
		\coordinate (A) at (0,\R);
		\coordinate (B) at (0,-\R);
		% OM = 2R, goc vuong tai A: MA^2 = OM^2 - OA^2 = 4R^2 - R^2 = 3R^2 => MA = R*sqrt(3)
		% OA thang dung len tren (0, R). M phai co OM = 2R va MA vuong goc OA => M nam tren duong thang y = R
		% x_M = -R*sqrt(3)
		\coordinate (M) at ({- \R*sqrt(3)}, \R);
		% Tiep diem C: O, C, M co MC la tiep tuyen. Goc AOM: cos(AOM) = OA/OM = 1/2 => AOM = 60 do.
		% Vay goc COM = 60 do nua, tuc la C o vi tri goc 150 do (hoac 90 + 2*60 = 210 do).
		% MA la tiep tuyen tai A(0, R) => M(-R*sqrt(3), R). Goc cua OA la 90 do.
		% Vector OM co goc la 150 do. Goc AOM = 60 do. Vay C co goc 150 + 60 = 210 do.
		\coordinate (C) at ({\R*cos(210)}, {\R*sin(210)});
		\coordinate (H) at (intersection of O--M and A--C);
		% MB cat duong tron tai D. B(0, -R).
		\coordinate (D) at (intersection of M--B and O--C); % tam thoi
		% Tinh toa do D chinh xac: giao diem MB va (O)
		\path [name path=circle] (O) circle (\R);
		\path [name path=MB] (M) -- (B);
		% Ta biet D la giao diem khac B cua duong thang MB voi (O).
		% MB di qua M(-sqrt(3)*R, R) va B(0, -R).
		% Vector MB = (sqrt(3)*R, -2R). He so goc: k = -2/sqrt(3).
		% Khoang cach tu O den MB: d = |0 - 0 - 2R*sqrt(3)| / sqrt(3 + 4) = 2*sqrt(3)/sqrt(7) * R ...
		\coordinate (D) at (-1.77,-2.17); % xap xi tren duong tron
		\coordinate (K) at ($(B)!0.5!(D)$);
		\coordinate (E) at (intersection of A--C and O--K);

		\draw[thick] (A) -- (B);
		\draw[thick] (M) -- (A);
		\draw[thick] (M) -- (C);
		\draw[thick] (M) -- (B);
		\draw[thick] (O) -- (M);
		\draw[thick] (A) -- (C);
		\draw[thick] (A) -- (D);
		\draw[dashed] (O) -- (D);
		\draw[dashed] (H) -- (D);
		\draw[dashed] (O) -- (E);
		\draw[dashed] (E) -- (B);

		\fill (O) circle (1.5pt) node[right] {$O$};
		\fill (A) circle (1.5pt) node[above] {$A$};
		\fill (B) circle (1.5pt) node[below] {$B$};
		\fill (M) circle (1.5pt) node[above left] {$M$};
		\fill (C) circle (1.5pt) node[below left] {$C$};
		\fill (H) circle (1.5pt) node[above right] {$H$};
		\fill (D) circle (1.5pt) node[below left] {$D$};
		\fill (K) circle (1.5pt) node[below] {$K$};
		\fill (E) circle (1.5pt) node[below] {$E$};
	\end{tikzpicture}
\end{center}
\begin{enumerate}[a)]
	\item Ta có điểm $D$ thuộc đường tròn đường kính $AB$, do đó $\widehat{ADB} = 90^\circ$ (góc nội tiếp chắn nửa đường tròn).\\
	Suy ra tam giác $ABD$ vuông tại $D$.\\
	Vì $MA$ và $MC$ là hai tiếp tuyến của $(O)$ cắt nhau tại $M$ nên $MA = MC$.\\
	Lại có $OA = OC = R$, suy ra $OM$ là đường trung trực của đoạn thẳng $AC$.\\
	Do đó $OM \perp AC$ tại $H$.
	\item Tam giác $OAM$ vuông tại $A$ có $AH$ là đường cao ($AH \perp OM$), áp dụng hệ thức lượng ta có:
	$$OA^2 = OH \cdot OM.$$
	Mà $OA = OD = R$ nên $OD^2 = OH \cdot OM \implies \dfrac{OD}{OH} = \dfrac{OM}{OD}$.\\
	Xét $\Delta ODH$ và $\Delta OMD$ có:
	\begin{itemize}
		\item Góc $\widehat{DOM}$ chung;
		\item $\dfrac{OD}{OH} = \dfrac{OM}{OD}$.
	\end{itemize}
	Do đó $\Delta ODH \backsim \Delta OMD$ (c-g-c), suy ra $\widehat{ODH} = \widehat{OMD}$.\\
	Mặt khác, trong tam giác vuông $OAM$, ta có $\widehat{OMA} = 90^\circ - \widehat{AOH} = \widehat{OAH} = \widehat{OAC}$.\\
	Mà tam giác $OAD$ cân tại $O$ ($OA = OD = R$) nên $\widehat{OAD} = \widehat{ODA}$.\\
	Suy ra $\widehat{ODH} = \widehat{DAC}$.
	\item Xét tam giác $MAB$ vuông tại $A$ ($MA$ là tiếp tuyến) có $AD$ là đường cao ($AD \perp MB$), ta có:
	$$\cos \widehat{AMB} = \dfrac{MA}{MB} = \dfrac{R\sqrt{3}}{\sqrt{MA^2 + AB^2}} = \dfrac{R\sqrt{3}}{\sqrt{3R^2 + 4R^2}} = \dfrac{\sqrt{3}}{\sqrt{7}}.$$
	Vì $K$ là trung điểm dây cung $BD$ nên $OK \perp BD$ tại $K$.\\
	Do đó tam giác $OKB$ vuông tại $K$, suy ra $\widehat{KOB} = 90^\circ - \widehat{OBK} = \widehat{BMA}$.\\
	Ta có $OM \perp AC$ tại $H$ và tam giác $OAC$ cân tại $O$ có $OH$ là phân giác nên $\widehat{AOH} = 60^\circ$.\\
	Trong tam giác vuông $OAE$, góc $\widehat{AOE} = 90^\circ$ hay $OE \perp AB$ tại $O$.\\
	Tam giác $OEB$ vuông tại $O$, có $OB = R$, $OE = OA \cdot \tan \widehat{OAE} = R \cdot \tan 30^\circ = \dfrac{R\sqrt{3}}{3}$.\\
	Diện tích tam giác vuông $OEB$ là:
	$$S_{OEB} = \dfrac{1}{2} OE \cdot OB = \dfrac{1}{2} \cdot \dfrac{R\sqrt{3}}{3} \cdot R = \dfrac{R^2\sqrt{3}}{6}.$$
\end{enumerate}
}
\end{bt}
